\documentclass[unicode,11pt,a4paper,oneside,numbers=endperiod,openany]{scrartcl}

\usepackage{ifthen}
\usepackage[utf8]{inputenc}
\usepackage{graphics}
\usepackage{graphicx}
\usepackage{hyperref}

\pagestyle{plain}
\voffset -5mm
\oddsidemargin  0mm
\evensidemargin -11mm
\marginparwidth 2cm
\marginparsep 0pt
\topmargin 0mm
\headheight 0pt
\headsep 0pt
\topskip 0pt        
\textheight 255mm
\textwidth 165mm

\newcommand{\duedate} {}
\newcommand{\setduedate}[1]{%
\renewcommand\duedate {See iCorsi for due date}}
\newcommand\isassignment {false}
\newcommand{\setassignment}{\renewcommand\isassignment {true}}
\newcommand{\ifassignment}[1]{\ifthenelse{\boolean{\isassignment}}{#1}{}}
\newcommand{\ifnotassignment}[1]{\ifthenelse{\boolean{\isassignment}}{}{#1}}

\newcommand{\assignmentpolicy}{
\begin{table}[h]
\begin{center}
\scalebox{0.8} {%
\begin{tabular}{|p{0.02cm}p{16cm}|}
\hline
&\\
\multicolumn{2}{|c|}{\Large\textbf{HPC Lab ---  Submission Instructions}}\\
\multicolumn{2}{|c|}{\large\textbf{(Please, notice that following instructions are mandatory: }}\\
\multicolumn{2}{|c|}{\large\textbf{submissions that don't comply with, won't be considered)}}\\
&\\
\textbullet & Assignments must be submitted to \href{https://www.icorsi.ch}{iCorsi} (i.e. in electronic format).\\
\textbullet & Provide source files (e.g. C/C++ files, Matlab). 
If you are using libraries, please add them in the file. Sources must be organized in directories called:\\
\multicolumn{2}{|c|}{\textit{Project\_number\_lastname\_firstname}}\\
& and  the  file must be called:\\
\multicolumn{2}{|c|}{\textit{project\_number\_lastname\_firstname.zip}}\\
\multicolumn{2}{|c|}{\textit{project\_number\_lastname\_firstname.pdf}}\\
\textbullet &  The TAs will grade your project by reviewing your project write-up, and looking at the implementation 
                 you attempted, and benchmarking your code's performance.\\

\textbullet & You are allowed to discuss all questions with anyone you like; however: (i) your submission must list anyone you discussed problems with and (ii) you must write up your submission independently.\\
\hline
\end{tabular}
}
\end{center}
\end{table}
}
\newcommand{\punkte}[1]{\hspace{1ex}\emph{\mdseries\hfill(#1~\ifcase#1{Points}\or{Points}\else{Points}\fi)}}


\newcommand\serieheader[6]{
\thispagestyle{empty}%
\begin{flushleft}
\includegraphics[width=0.4\textwidth]{usi_inf.png}
\end{flushleft}
  \noindent%
  {\large\ignorespaces{\textbf{#1}}\hspace{\fill}\ignorespaces{ \textbf{#2}}}\\ \\%
  {\large\ignorespaces #3 \hspace{\fill}\ignorespaces #4}\\
  \noindent%
  \bigskip
  \hrule\par\bigskip\noindent%
  \bigskip {\ignorespaces {\Large{\textbf{#5}}}
  \hspace{\fill}\ignorespaces \large \ifthenelse{\boolean{\isassignment}}{\duedate}{#6}}
  \hrule\par\bigskip\noindent%  \linebreak
 }

\makeatletter
\def\enumerateMod{\ifnum \@enumdepth >3 \@toodeep\else
      \advance\@enumdepth \@ne
      \edef\@enumctr{enum\romannumeral\the\@enumdepth}\list
      {\csname label\@enumctr\endcsname}{\usecounter
        {\@enumctr}%%%? the following differs from "enumerate"
	\topsep0pt%
	\partopsep0pt%
	\itemsep0pt%
	\def\makelabel##1{\hss\llap{##1}}}\fi}
\let\endenumerateMod =\endlist
\makeatother




\usepackage{textcomp}





\usepackage{listings}
\begin{document}


\setassignment

\serieheader{High-Performance Computing Lab}{Institute of Computing}{Student: Paolo Deidda}{Discussed with: FULL NAME}{Solution for Project 1}{}
\newline

\assignmentpolicy
In this project you will practice memory access optimization, performance-oriented programming, and OpenMP parallelizaton 
on the Rosa Cluster .  

\section{Rosa Warm-Up \punkte{5}}

\begin{enumerate}
    \item % 1
    The \textbf{module system} is a utility that allows the user to dynamically manage their software environment on the Rosa HPC cluster and to load different compilers, libraries and applications in order to modify environment variagles (like PATH, LD LIBRARY PATH, MANPATH, etc.) without creating conflicts between different software versions or dependencies.    

    As riported in the USI resource page the module system provides several commands to manage the environment.

    \begin{itemize}
        \item \texttt{module avail} -- lists all available modules (on the current system)
        \item \texttt{module list} -- lists all currently loaded modules
        \item \texttt{module show} -- display information about
        \item \texttt{module load} -- loads module
        \item \texttt{module switch} -- unloads, loads
        \item \texttt{module rm} -- unloads module
        \item \texttt{module purge} -- unloads all loaded modules
    \end{itemize}

    \item % 2
    The \textbf{Slurm} (Simple Linux Utility for Resource Management) is a job scheduler for Linux clusters.
    Main features of Slurm are:
    \begin{itemize}
        \item Job scheduling and resource management
        \item Framework for starting, executing, and monitoring work (jobs) on a set of allocated nodes
        \item Queuing management to handle multiple users and jobs
    \end{itemize}

    The two main components are:
    \begin{itemize}
        \item \textit{sulurmd}: the deamnon that runs on each compute node responsible for launching, monitoring, and terminating jobs
        \item \textit{slurmctld}: the central management daemon that manages job queues and allocates resources
    \end{itemize}

    Main commands:
    \begin{itemize}
        \item \texttt{srun}: submit a job for execution
        \item \texttt{sbatch}: submit a batch job
        \item \texttt{squeue}: view the status of jobs in the queue
        \item \texttt{scancel}: cancel a job
        \item \texttt{salloc}: allocate resources for an interactive job
    \end{itemize}

    \item % 3
    The program is at \texttt{src/1-Rosa-warm-up/hello\_world.c} and the code is the following:
    % import file content
    \lstinputlisting[language=C]{../src/1-Rosa-warm-up/hello_world.c}

    \item % 4
    We can see the output of the command sinfo here below: 
    \lstinputlisting[language=bash]{../src/1-Rosa-warm-up/sinfo_output.txt}
    
    
\end{enumerate}




\section{Performance Characteristics \punkte{30}}


\section{Optimize Square Matrix-Matrix Multiplication  \punkte{50}}


\section{Quality of the Report  \punkte{15}}


\end{document}
